\documentclass{beamer}

%% Tema y configuraciones
\usetheme{Boadilla}
\usecolortheme{default}
\setbeamertemplate{navigation symbols}{}
\usepackage{xcolor}
\usepackage{tikz}
\usetikzlibrary{arrows.meta, positioning}
\hypersetup{
    colorlinks=true,
    linkcolor=.,      % color de enlaces internos (índices, refs)
    urlcolor=blue,    % color de URLs
    citecolor=.,      % color de referencias bibliográficas
}

%% Helpers (macros y notación)
\newcommand{\E}{\mathbb{E}}
\newcommand{\Var}{\mathrm{Var}}
\newcommand{\Cov}{\mathrm{Cov}}
\newcommand{\1}{\mathbb{I}}
\newcommand{\MSE}{\mathrm{MSE}}
\DeclareMathOperator*{\argmin}{arg\,min}

%% Encabezado
\title[BDML]{Big Data y Machine Learning \\ para Economía Aplicada}
\subtitle{Week 08: Redes Neuronales y Visión por Computadora}
\author{Eduard F. Martinez Gonzalez, Ph.D.}
\institute[]{Departamento de Economía, Universidad Icesi}
\date{\today}

%%=================================================
%% Begin document
\begin{document}

%%=================================================
%% Primer slide
\begin{frame}
\titlepage
\end{frame}

%%=================================================
\begin{frame}{Previously\ldots}
\small
La frase con la que cerró la sesión pasada: \textbf{representar es comprimir con sentido} --- y PCA lo hacía \emph{linealmente}.
\pause

\vspace{0.15cm}
Hoy convergen todas las piezas que hemos ido sembrando:
\begin{itemize}\itemsep4pt
  \item El \textbf{gradiente} como motor de aprendizaje (\emph{boosting}, S4; verosimilitud logística, S5).
  \item La \textbf{regularización} como precio de la flexibilidad (S2--S3).
  \item La \textbf{sigmoide} y la \emph{cross-entropy} (S5).
  \item Las \textbf{representaciones} aprendidas (S7).
\end{itemize}
\pause

\vspace{0.15cm}
\begin{block}{Hoy: el aproximador universal}
\begin{enumerate}\itemsep2pt
  \item \textbf{Redes feedforward}: qué son, por qué aproximan cualquier función, y cómo se entrenan (\textbf{retropropagación}, derivada en tablero: va al Taller 2).
  \item \textbf{Redes convolucionales}: la arquitectura que lee imágenes.
  \item \textbf{Economía desde el cielo}: predecir pobreza con imágenes satelitales.
\end{enumerate}
\end{block}
\end{frame}


%%#################################################
%%#################################################
%% PARTE 1: DEL PERCEPTRÓN A LA RED
%%#################################################
%%#################################################
\section{Del perceptrón a la red}
\begin{frame}[noframenumbering]
\tableofcontents[currentsection]
\end{frame}

%%=================================================
\begin{frame}{La neurona: una regresión con actitud}
\small
La unidad básica no tiene ningún misterio:
\[ a \;=\; g\big( w'x + b \big) \]
combinación lineal de las entradas $+$ una \textbf{función de activación} no lineal $g$.

\vspace{0.15cm}
\begin{center}
\begin{tikzpicture}[scale=0.78, every node/.style={font=\scriptsize}]
  \node[circle, draw] (x1) at (0,1.5) {$x_1$};
  \node[circle, draw] (x2) at (0,0.5) {$x_2$};
  \node[circle, draw] (x3) at (0,-0.5) {$x_3$};
  \node[circle, draw, fill=blue!10, minimum size=1.0cm] (n) at (2.6,0.5) {$\Sigma \,\big|\, g$};
  \node (out) at (4.7,0.5) {$a = g(w'x + b)$};
  \draw[-{Stealth}] (x1) -- node[above, sloped]{$w_1$} (n);
  \draw[-{Stealth}] (x2) -- node[above]{$w_2$} (n);
  \draw[-{Stealth}] (x3) -- node[below, sloped]{$w_3$} (n);
  \draw[-{Stealth}, thick] (n) -- (out);
\end{tikzpicture}
\end{center}
\pause

\begin{itemize}\itemsep4pt
  \item Con $g$ sigmoide, una neurona \textbf{es} una regresión logística (S5). Nada nuevo aún.
  \item La novedad no es la neurona: es \textbf{componerlas}. \pause
  \item Sin la no linealidad no hay negocio: componer funciones lineales da\ldots\ otra función lineal. $g$ es lo que hace que apilar capas agregue capacidad.
\end{itemize}
\end{frame}

%%=================================================
\begin{frame}{La red feedforward: composición de capas}
\small
\begin{center}
\begin{tikzpicture}[scale=0.66, every node/.style={font=\scriptsize}]
  % input
  \foreach \i/\y in {1/1.0, 2/0, 3/-1.0} \node[circle, draw, minimum size=0.62cm] (i\i) at (0,\y) {$x_\i$};
  % hidden
  \foreach \i/\y in {1/1.5, 2/0.5, 3/-0.5, 4/-1.5} \node[circle, draw, fill=blue!10, minimum size=0.62cm] (h\i) at (2.6,\y) {};
  % output
  \foreach \i/\y in {1/0.5, 2/-0.5} \node[circle, draw, fill=red!10, minimum size=0.62cm] (o\i) at (5.2,\y) {};
  \foreach \i in {1,2,3} \foreach \j in {1,2,3,4} \draw[gray!60] (i\i) -- (h\j);
  \foreach \i in {1,2,3,4} \foreach \j in {1,2} \draw[gray!60] (h\i) -- (o\j);
  \node at (0,-2.15) {entrada $x$};
  \node at (2.6,-2.35) {capa oculta $a^{(1)}$};
  \node at (5.2,-2.15) {salida $\hat{y}$};
\end{tikzpicture}
\end{center}
\pause

\textbf{En notación matricial, capa por capa} (con $a^{(0)} = x$):
\[ a^{(l)} = g\big( W^{(l)} a^{(l-1)} + b^{(l)} \big), \quad l = 1, \ldots, L-1; \qquad \hat{y} = h\big( W^{(L)} a^{(L-1)} + b^{(L)} \big) \]
\pause

\begin{itemize}\itemsep3pt
  \item La red \textbf{es} una composición: $\hat{y} = f_L \circ \cdots \circ f_1(x)$. Profundidad $=$ cuántas funciones componemos.
  \item Los ``parámetros'' son todos los $W, b$: fácilmente miles o millones --- flexibilidad brutal que exigirá regularización brutal (lo sabíamos desde la S1).
  \item Salida $h$ según la tarea: identidad (regresión), sigmoide (binaria), \textbf{softmax} ($K$ clases).
\end{itemize}
\end{frame}

%%=================================================
\begin{frame}{¿Por qué tanta fe? Aproximación universal y profundidad}
\small
\textbf{Teorema de aproximación universal} (Cybenko, 1989; Hornik, 1991): una red con \textbf{una} capa oculta suficientemente ancha aproxima cualquier función continua tan bien como se quiera.
\pause

\begin{itemize}\itemsep5pt
  \item Suena definitivo, pero es una garantía de \emph{existencia}: no dice cuántas neuronas (pueden ser astronómicas) ni cómo encontrar los pesos. \pause
  \item \textbf{¿Por qué profundo y no ancho, entonces?} La composición \textbf{reutiliza}: las capas tempranas aprenden piezas (bordes, patrones básicos) que las tardías \emph{combinan}. Funciones con estructura composicional se representan con muchos menos parámetros en profundidad que en anchura. \pause
  \item Y la lectura que conecta con la sesión 7: \textbf{cada capa oculta es una representación aprendida} de los datos --- de cruda a abstracta. La red no solo predice: fabrica las variables con las que predice.
\end{itemize}
\pause

\begin{block}{El contrato de la flexibilidad}
Aproximarlo todo $\Rightarrow$ sobreajustarlo todo. El teorema garantiza capacidad; la generalización se compra con datos y regularización --- la curva U jamás se fue.
\end{block}
\end{frame}

%%=================================================
\begin{frame}{Funciones de activación}
\small
\begin{center}
\begin{tikzpicture}[scale=0.6]
% sigmoide
\begin{scope}[shift={(0,0)}]
  \draw[->] (-2.6,0) -- (2.6,0); \draw[->] (0,-0.3) -- (0,2.5);
  \draw[dashed, gray] (-2.4,2.0) -- (2.4,2.0);
  \draw[blue, very thick] plot[smooth] coordinates {(-2.5,0.11) (-1.5,0.36) (-0.7,0.66) (0,1.0) (0.7,1.34) (1.5,1.64) (2.5,1.89)};
  \node[font=\footnotesize] at (0,-0.85) {sigmoide};
\end{scope}
% tanh
\begin{scope}[shift={(6.4,1.0)}]
  \draw[->] (-2.6,0) -- (2.6,0); \draw[->] (0,-1.4) -- (0,1.5);
  \draw[red, very thick] plot[smooth] coordinates {(-2.5,-0.99) (-1.5,-0.91) (-1,-0.76) (-0.5,-0.46) (0,0) (0.5,0.46) (1,0.76) (1.5,0.91) (2.5,0.99)};
  \node[font=\footnotesize] at (0,-1.85) {tanh};
\end{scope}
% relu
\begin{scope}[shift={(12.8,0)}]
  \draw[->] (-2.6,0) -- (2.6,0); \draw[->] (0,-0.3) -- (0,2.5);
  \draw[teal, very thick] (-2.4,0) -- (0,0) -- (2.2,2.2);
  \node[font=\footnotesize] at (0,-0.85) {ReLU $= \max(0, z)$};
\end{scope}
\end{tikzpicture}
\end{center}
\pause

\begin{itemize}\itemsep5pt
  \item \textbf{Sigmoide/tanh}: históricas y elegantes, pero se \textbf{saturan}: en las colas su derivada es $\approx 0$, y con muchas capas el gradiente se \emph{desvanece} multiplicándose por casi-ceros.
  \item \textbf{ReLU}: tosca (¡ni siquiera es diferenciable en 0!) pero no se satura en la mitad positiva, es baratísima, y en la práctica entrena mucho mejor: el \textbf{estándar} en capas ocultas.
  \item \textbf{Softmax} en la salida multiclase: convierte $K$ puntajes en probabilidades ($e^{z_k}/\sum_j e^{z_j}$) --- la generalización de la logística que ya conocen.
\end{itemize}
\end{frame}


%%#################################################
%%#################################################
%% PARTE 2: ENTRENAMIENTO
%%#################################################
%%#################################################
\section{Entrenamiento: SGD y retropropagación}
\begin{frame}[noframenumbering]
\tableofcontents[currentsection]
\end{frame}

%%=================================================
\begin{frame}{La pérdida, y una mala noticia asumida con calma}
\small
\textbf{Nada nuevo en la pérdida:}
\begin{itemize}\itemsep3pt
  \item Regresión: $L = \frac{1}{n}\sum_i (y_i - \hat{y}_i)^2$ (el MSE de siempre).
  \item Clasificación: \emph{cross-entropy} $= -\frac{1}{n}\sum_i \sum_k y_{ik} \log \hat{p}_{ik}$ --- que es, literalmente, la \textbf{máxima verosimilitud} de la sesión 5 con signo menos.
\end{itemize}
\pause

\vspace{0.2cm}
\textbf{La mala noticia:} como función de los pesos $(W, b)$, la pérdida \textbf{no es convexa}: paisaje lleno de valles, mesetas y puntos de silla. Adiós garantías globales tipo Lasso.
\pause

\vspace{0.15cm}
\textbf{La sorpresa empírica que sostiene todo el campo:} el descenso de gradiente \emph{estocástico}, sin garantía teórica de óptimo global, encuentra \textbf{sistemáticamente} soluciones que generalizan bien en redes grandes. Los mínimos ``suficientemente buenos'' abundan.
\pause

\vspace{0.15cm}
\begin{block}{Actitud del curso}
Igual que con $k$-means (S7): óptimo local $+$ buen protocolo (inicializaciones, monitoreo, regularización) $=$ herramienta utilizable. La perfección teórica no es requisito para la utilidad práctica.
\end{block}
\end{frame}

%%=================================================
\begin{frame}{Descenso de gradiente estocástico (SGD)}
\small
\textbf{La regla de siempre}, ahora sobre millones de parámetros $\theta = (W, b)$:
\[ \theta \;\leftarrow\; \theta \;-\; \eta \, \nabla_\theta L(\theta) \]
\pause

\textbf{Las tres variantes:}
\begin{itemize}\itemsep4pt
  \item \textbf{Batch}: gradiente con \emph{todos} los datos por paso --- exacto y carísimo.
  \item \textbf{Estocástico puro}: una observación por paso --- barato y ruidosísimo.
  \item \textbf{Mini-batch} (el estándar): 32--256 observaciones por paso: gradiente ruidoso pero decente, y el ruido hasta ayuda a escapar de valles malos.
\end{itemize}
\pause

\begin{itemize}\itemsep4pt
  \item $\eta$ (\emph{learning rate}): \textbf{la} perilla delicada --- grande diverge, chica no llega. Se afina con \emph{schedules} o con \textbf{Adam} (pasos adaptativos por parámetro; el default moderno).
  \item Una \textbf{época} $=$ una pasada por los datos; se entrena por épocas monitoreando validación. \pause
\end{itemize}

\begin{block}{Deja vu del boosting}
En la S4 hicimos descenso de gradiente \emph{en el espacio de funciones} (un árbol por paso). Aquí es \emph{en el espacio de parámetros} (un mini-batch por paso). Mismo motor, distinto chasis.
\end{block}
\end{frame}

%%=================================================
\begin{frame}{Retropropagación (I): la idea}
\small
\textbf{El problema computacional:} $\nabla_\theta L$ tiene millones de componentes. ¿Cómo calcularlos sin morir?
\pause

\vspace{0.2cm}
\textbf{La respuesta: la regla de la cadena, administrada con inteligencia.}
\begin{itemize}\itemsep5pt
  \item La red es una composición $\Rightarrow$ la derivada de la pérdida respecto a \emph{cualquier} peso es un producto de derivadas locales a lo largo del camino hasta la salida.
  \item La observación clave: esos productos \textbf{comparten factores}. Computando \emph{de la salida hacia atrás}, cada factor se calcula \textbf{una sola vez} y se reutiliza. \pause
  \item Resultado: el gradiente completo cuesta lo mismo que \textbf{una pasada hacia adelante $+$ una hacia atrás} --- no una pasada por parámetro. Esa economía es lo que hace entrenables las redes gigantes.
\end{itemize}
\pause

\vspace{0.15cm}
\begin{itemize}\itemsep3pt
  \item El objeto que viaja hacia atrás: el ``error'' de cada capa, $\delta^{(l)} = \partial L / \partial z^{(l)}$.
  \item En la práctica nadie deriva a mano: \textbf{diferenciación automática} (Keras/PyTorch). Pero ustedes sí lo harán una vez --- ahora mismo --- porque es material del \textbf{Taller 2}.
\end{itemize}
\end{frame}

%%=================================================
\begin{frame}{Retropropagación (II): derivación con una capa oculta}
\small
\textbf{Setup} (regresión, pérdida cuadrática de una observación):
\[ z = W^{(1)} x + b^{(1)}, \quad h = g(z), \quad \hat{y} = w^{(2)\prime} h + b^{(2)}, \quad L = \tfrac{1}{2}\big(\hat{y} - y\big)^2 \]
\pause
\textbf{1. Capa de salida} (regla de la cadena, primer eslabón):
\[ \delta^{(2)} \;=\; \frac{\partial L}{\partial \hat{y}} \;=\; \hat{y} - y \qquad \text{(¡el residuo, otra vez!)} \]
\pause
\textbf{2. Gradientes de la capa de salida:}
\[ \frac{\partial L}{\partial w^{(2)}} = \delta^{(2)}\, h, \qquad \frac{\partial L}{\partial b^{(2)}} = \delta^{(2)} \]
\pause
\textbf{3. Propagar el error hacia atrás} (por cada neurona oculta, la cadena pasa por $\hat{y}$ y por $g$):
\[ \delta^{(1)} \;=\; \big( w^{(2)} \, \delta^{(2)} \big) \odot g'(z) \]
\pause
\textbf{4. Gradientes de la capa oculta:}
\[ \frac{\partial L}{\partial W^{(1)}} = \delta^{(1)} x', \qquad \frac{\partial L}{\partial b^{(1)}} = \delta^{(1)} \]
\pause
\begin{block}{Para el Taller 2}
``Un paso de retropropagación en una red simple'' $=$ estos cuatro pasos, con números. El patrón general: $\delta$ de una capa $=$ ($\delta$ de la siguiente $\times$ pesos) $\odot$ derivada de la activación.
\end{block}
\end{frame}

%%=================================================
\begin{frame}{Regularizar redes: el kit completo}
\small
Millones de parámetros $\Rightarrow$ el sobreajuste no es riesgo, es \textbf{default}. El arsenal (todo les va a sonar):
\pause
\begin{itemize}\itemsep5pt
  \item \textbf{Weight decay} ($L_2$): penalizar $\sum \|W\|^2$ --- Ridge (S2) con otro nombre y la misma alma.
  \item \textbf{Early stopping}: monitorear la pérdida de \emph{validación} y parar cuando sube --- es \textbf{elegir el mínimo de la curva U en vivo} (S1--S2), con las épocas como eje de complejidad.
  \item \textbf{Dropout} (Srivastava et al., 2014): en cada paso, apagar al azar una fracción de neuronas $\Rightarrow$ la red no puede depender de nadie en particular; en el fondo, un \textbf{ensamble implícito} de sub-redes (¡el espíritu del Random Forest, S4!).
  \item \textbf{Batch normalization}: estandarizar activaciones dentro de la red (nuestra vieja obsesión por estandarizar, ahora capa por capa); estabiliza y acelera.
  \item \textbf{Aumentar datos} (\emph{augmentation}): en imágenes, rotar/recortar/reflejar --- fabricar variación invariante. El regularizador más honesto: más datos.
\end{itemize}
\pause

\begin{block}{Cuándo NO usar redes}
Datos tabulares medianos: XGBoost/RF suelen ganar con una fracción del esfuerzo (S4). Las redes pagan cuando hay \textbf{estructura} (imágenes, texto, secuencias) y \textbf{escala}.
\end{block}
\end{frame}


%%#################################################
%%#################################################
%% PARTE 3: CNN — REDES QUE VEN
%%#################################################
%%#################################################
\section{Redes convolucionales: visión}
\begin{frame}[noframenumbering]
\tableofcontents[currentsection]
\end{frame}

%%=================================================
\begin{frame}{Por qué una red densa no puede ``ver''}
\small
Una imagen modesta de $200 \times 200$ píxeles a color $= 120.000$ entradas.
\pause
\begin{itemize}\itemsep5pt
  \item Primera capa densa de 1.000 neuronas $\Rightarrow$ \textbf{120 millones} de pesos\ldots\ en una capa. Sobreajuste garantizado y cómputo absurdo. \pause
  \item Peor: la red densa \textbf{ignora la estructura espacial}. Trata el píxel (1,1) y el (200,200) como si su relación fuera igual a la de dos vecinos --- y si el gato se corre 10 píxeles, para la red es una entrada completamente nueva. \pause
\end{itemize}

\vspace{0.15cm}
\textbf{Las dos ideas que arreglan todo (LeCun, 1989):}
\begin{enumerate}\itemsep4pt
  \item \textbf{Localidad}: los patrones visuales son locales (un borde vive en una ventanita, no en toda la imagen).
  \item \textbf{Invarianza a traslación}: un borde es un borde \emph{esté donde esté} $\Rightarrow$ usar \textbf{los mismos pesos} en toda la imagen.
\end{enumerate}
\pause
Juntas se llaman \textbf{convolución} --- y convierten los 120 millones de pesos en unas decenas por filtro.
\end{frame}

%%=================================================
\begin{frame}{La convolución: un filtro que recorre la imagen}
\small
\begin{center}
\begin{tikzpicture}[scale=0.78]
  % input grid 6x6
  \fill[blue!18] (0.4,0.8) rectangle (1.6,2.0);
  \draw[step=0.4, gray] (0,0) grid (2.4,2.4);
  \node[font=\scriptsize] at (1.2,-0.4) {imagen (p.\,ej.\ $6\times6$)};
  % kernel
  \fill[orange!25] (3.4,1.0) rectangle (4.6,2.2);
  \draw[step=0.4, gray] (3.4,1.0) grid (4.6,2.2);
  \node[font=\scriptsize, align=center] at (4.0,0.35) {filtro $3\times3$ \\ (pesos compartidos)};
  % arrow
  \draw[-{Stealth}, thick] (5.0,1.5) -- (5.9,1.5);
  % output grid 4x4
  \fill[blue!35] (6.4,1.2) rectangle (6.8,1.6);
  \draw[step=0.4, gray] (6.0,0.4) grid (7.6,2.0);
  \node[font=\scriptsize, align=center] at (6.8,-0.25) {mapa de \emph{features} \\ ($4\times4$)};
\end{tikzpicture}
\end{center}
\pause

\begin{itemize}\itemsep4pt
  \item El filtro (una matriz chica de pesos) se \textbf{desliza} por la imagen; en cada posición: producto punto ventana--filtro $\rightarrow$ un número del mapa de salida.
  \item Cada filtro aprende a \textbf{detectar un patrón} (un borde vertical, una esquina, una textura); su mapa dice \emph{dónde} apareció.
  \item Una capa convolucional $=$ varios filtros en paralelo; los pesos a aprender son los filtros --- decenas, no millones. \pause
  \item \textbf{Pooling} (submuestreo, p.\,ej.\ el máximo de cada $2\times2$): reduce resolución, gana robustez a desplazamientos pequeños y abarata lo que sigue.
\end{itemize}
\end{frame}

%%=================================================
\begin{frame}{La arquitectura CNN: jerarquía de representaciones}
\small
\textbf{El patrón que se repite:} $\big[$ convolución $\rightarrow$ ReLU $\rightarrow$ pooling $\big] \times$ varias veces $\rightarrow$ capas densas $\rightarrow$ softmax.
\pause

\vspace{0.15cm}
\textbf{Lo que aprende cada nivel} (verificado mirando dentro de redes entrenadas):
\begin{itemize}\itemsep3pt
  \item Capas tempranas: \textbf{bordes} y gradientes de color.
  \item Capas medias: \textbf{texturas y partes} (ventanas, ruedas, techos, esquinas).
  \item Capas tardías: \textbf{objetos y escenas} (una casa, una vía, un cultivo).
\end{itemize}
La \textbf{composición} de la Parte 1, hecha visible: representaciones de lo simple a lo abstracto. \pause

\vspace{0.15cm}
\textbf{El momento histórico:} AlexNet gana ImageNet (2012) por paliza --- profundidad $+$ ReLU $+$ dropout $+$ GPUs. Desde entonces, la visión por computadora \textbf{es} aprendizaje profundo.
\pause

\vspace{0.1cm}
\begin{block}{Transfer learning: el truco que nos habilita a nosotros}
Una CNN entrenada en millones de imágenes genéricas ya sabe ver bordes, texturas y partes --- eso es \textbf{universal}. Receta: tomar la red preentrenada, \textbf{congelar} las capas tempranas y reentrenar solo las finales con \emph{nuestros} datos (miles, no millones). Así se hace visión con presupuesto de economista.
\end{block}
\end{frame}


%%#################################################
%%#################################################
%% PARTE 4: ECONOMÍA DESDE EL CIELO
%%#################################################
%%#################################################
\section{Economía desde el cielo}
\begin{frame}[noframenumbering]
\tableofcontents[currentsection]
\end{frame}

%%=================================================
\begin{frame}{Jean et al.\ (2016, Science): pobreza vista desde satélites}
\small
\textbf{El problema:} en buena parte del mundo en desarrollo, las encuestas de hogares son escasas, caras y poco frecuentes --- justo donde más se necesita medir pobreza.
\pause

\vspace{0.15cm}
\textbf{La solución, paso a paso} (transfer learning doblemente inteligente):
\begin{enumerate}\itemsep4pt
  \item Partir de una CNN \textbf{preentrenada} en imágenes genéricas (ImageNet).
  \item Reentrenarla para predecir \textbf{luces nocturnas} a partir de imágenes satelitales \textbf{diurnas} --- las luces son un proxy conocido de actividad económica, y hay datos de sobra (no se necesita ni una encuesta para este paso).
  \item En ese proceso, la red aprende sola a reconocer techos, vías, cultivos, urbanización\ldots
  \item Usar esas \emph{features} para predecir el \textbf{consumo medido en las encuestas} disponibles.
\end{enumerate}
\pause

\textbf{Resultado:} explica entre \textbf{37\% y 55\%} de la variación del consumo promedio local en cinco países africanos --- con datos esencialmente \textbf{gratis y actualizables}.
\pause

\begin{block}{Por qué es la aplicación perfecta del curso}
Transfer learning (hoy) $+$ proxies y validación \emph{downstream} (S7) $+$ evaluación out-of-sample honesta (S1--S2), al servicio de un problema de medición económica real.
\end{block}
\end{frame}

%%=================================================
\begin{frame}{El ecosistema: luces, píxeles y política pública}
\small
\textbf{La familia de aplicaciones que abre esta caja:}
\begin{itemize}\itemsep5pt
  \item \textbf{Luces nocturnas como PIB local} (Henderson, Storeygard \& Weil, 2012, AER): medir crecimiento donde las cuentas nacionales son débiles o dudosas; hoy, estándar para PIB municipal y conflicto. \pause
  \item \textbf{Deforestación y uso del suelo}: monitoreo casi en tiempo real con clasificación de píxeles --- la base de la fiscalización ambiental moderna (Amazonía incluida). \pause
  \item \textbf{Urbanización y vivienda}: detectar construcciones, medir expansión informal, actualizar catastros --- con obvia resonancia colombiana. \pause
  \item \textbf{Focalización}: mapas de pobreza de alta resolución para asignar programas donde las encuestas no llegan.
\end{itemize}
\pause

\vspace{0.1cm}
\begin{block}{La advertencia de medición (siempre la hay)}
Lo predicho es un \textbf{proxy}: hereda los sesgos de su entrenamiento (¿qué pobreza \emph{no} se ve desde arriba?) y se degrada fuera del dominio (otro país, otra década). Validar contra terreno --- la lección de la S7 --- no es opcional.
\end{block}
\end{frame}

%%=================================================
\begin{frame}{Recap}
\small
\begin{enumerate}\itemsep6pt
  \item Una neurona es una logística; la potencia está en la \textbf{composición}: capas $=$ representaciones aprendidas, de lo crudo a lo abstracto.
  \item \textbf{Aproximación universal} garantiza capacidad, no generalización: la curva U manda, y el kit es conocido --- weight decay (Ridge), early stopping (la U en vivo), dropout (ensamble implícito), batch norm, más datos.
  \item Entrenar $=$ \textbf{SGD con mini-batches} (Adam como default) sobre una pérdida no convexa que en la práctica se deja optimizar.
  \item \textbf{Retropropagación} $=$ regla de la cadena con reutilización: $\delta^{(2)} = \hat{y} - y$; $\delta^{(1)} = (w^{(2)}\delta^{(2)}) \odot g'(z)$; gradiente al costo de ida $+$ vuelta. (Taller 2.)
  \item \textbf{CNN}: localidad $+$ pesos compartidos $=$ ver con pocos parámetros; jerarquía bordes $\rightarrow$ partes $\rightarrow$ objetos; \textbf{transfer learning} para usarlas con datos de economista.
  \item \textbf{Satélites}: de luces nocturnas (Henderson et al.) a mapas de pobreza (Jean et al.): medición barata, actualizable --- y siempre un proxy que hay que validar.
\end{enumerate}
\end{frame}

%%=================================================
\begin{frame}
\vspace{2.0cm}
\Large{Preparación para la próxima clase\ldots}
\end{frame}

%%#################################################
%%#################################################
%% PARTE 5: PRÓXIMA CLASE
%%#################################################
%%#################################################

%%=================================================
\begin{frame}{Para aprovechar al máximo la próxima sesión}
\small
\textbf{Lecturas obligatorias de esta semana:}
\begin{itemize}\itemsep4pt
  \item Goodfellow, Bengio \& Courville (2016), \emph{Deep Learning}, cap.\ 6 (secs.\ 6.1--6.5). \\ \textcolor{gray}{Guía: hasta backprop; el cap.\ 7 (regularización) es el complemento natural.}
  \item Jean et al.\ (2016), \emph{Science}. \\ \textcolor{gray}{Guía: es corto; concéntrense en la figura 1 y el diseño del transfer learning.}
\end{itemize}

\textbf{Recomendadas:} Henderson, Storeygard \& Weil (2012); LeCun, Bengio \& Hinton (2015, \emph{Nature}, el survey).
\pause

\vspace{0.2cm}
\textbf{Próxima sesión --- la última de contenido:} \textbf{texto como datos y LLMs}. Del conteo de palabras a los \emph{embeddings}, la atención y los \emph{transformers}: entrenaremos un GPT de juguete en vivo y haremos \emph{fine-tuning} de un modelo abierto. \textbf{Logística:} traigan portátil con cuenta de Google activa (usaremos Colab con GPU gratuita); la sesión corre en \textbf{Python} con cuadernos preparados --- la excepción anunciada a la regla ``R primero''.
\pause

\vspace{0.2cm}
\textbf{Aplicación de esta semana en R:} red feedforward con \texttt{keras} para predecir consumo eléctrico municipal, y clasificación de imágenes con una CNN preentrenada.
\end{frame}

%%=================================================
%% Referencias
\begin{frame}{Referencias esenciales}
\footnotesize
\begin{itemize}\itemsep2pt
  \item Goodfellow, I., Bengio, Y., \& Courville, A. (2016). \emph{Deep Learning}. MIT Press. Caps.\ 6--9.
  \item Cybenko, G. (1989). Approximation by Superpositions of a Sigmoidal Function. \emph{Math.\ of Control, Signals and Systems}, 2, 303--314.
  \item LeCun, Y., Bottou, L., Bengio, Y., \& Haffner, P. (1998). Gradient-Based Learning Applied to Document Recognition. \emph{Proc.\ IEEE}, 86(11).
  \item Krizhevsky, A., Sutskever, I., \& Hinton, G. (2012). ImageNet Classification with Deep Convolutional Neural Networks. \emph{NeurIPS}.
  \item Srivastava, N., et al.\ (2014). Dropout: A Simple Way to Prevent Neural Networks from Overfitting. \emph{JMLR}, 15, 1929--1958.
  \item LeCun, Y., Bengio, Y., \& Hinton, G. (2015). Deep Learning. \emph{Nature}, 521, 436--444.
  \item Henderson, V., Storeygard, A., \& Weil, D. (2012). Measuring Economic Growth from Outer Space. \emph{AER}, 102(2), 994--1028.
  \item Jean, N., Burke, M., Xie, M., Davis, W., Lobell, D., \& Ermon, S. (2016). Combining Satellite Imagery and Machine Learning to Predict Poverty. \emph{Science}, 353(6301), 790--794.
  \item Dell, M. (2025). Deep Learning for Economists. \emph{Journal of Economic Literature}, 63(1), 5--58.
\end{itemize}
\normalsize
\end{frame}

%%=================================================
%% End document
\end{document}
